\chapter{Traceability between Functional Requirements and Competency Questions}
\label{ap:OntoKSD-FR-CQ}

This appendix presents the requirement-oriented view of the revised
OntoKSD competency-question collection. Table~\ref{tab:traceability-fr-cq}
complements the question-oriented mapping in
Table~\ref{tab:traceability-cq-fr}, using the same requirement definitions,
question families, and coverage limits.

The mapping refers to the revised question formulations and their
implemented retrieval scope. It identifies which aspects of a requirement
can be examined through the related queries and which require additional
checks. A mapping is not a statement of requirement satisfaction:
preparing a query, executing it without error, and obtaining the expected
answer provide different kinds of evidence. Requirement-level conclusions
must also consider the dataset, parameters, inference configuration,
expected results, and applicable representational checks.

The identifiers CQ1--CQ42 denote historical question families retained
after the qualification review. Their 58 operational files include
complementary views, equivalent retrievals, and one optional update.
Suffixes preserve these relationships; they do not identify additional
functional requirements or independent successful tests. The repository
catalog and identifier map document the corresponding files.

\section{Requirement-to-Question Mapping}

\begingroup
\small
\begin{longtable}{
    p{0.08\linewidth}
    p{0.31\linewidth}
    p{0.49\linewidth}
}
\caption{Functional requirements, related CQ families, and complementary verification needs.}
\label{tab:traceability-fr-cq} \\
\toprule
\textbf{FR} & \textbf{Related CQ families} & \textbf{Scope and complementary checks} \\
\midrule
\endfirsthead

\toprule
\textbf{FR} & \textbf{Related CQ families} & \textbf{Scope and complementary checks} \\
\midrule
\endhead

\bottomrule
\endfoot

FR1 & CQ1, CQ4, CQ5, CQ8 &
Retrieve specification identity, constituent organization, pattern, and
represented structure. Examine scope descriptions and completeness
against the source specification separately. \\

FR2 & CQ2, CQ3 &
Retrieve knowledge--level pairings and recorded justification verbs.
Verify the level's framework-specific typing and reference scheme in
the imported vocabulary, and distinguish required skill levels from
learner proficiency. \\

FR3 & CQ1 &
Retrieve dispositional constituents. Examine their meaning and relation
to expected action; inclusion in a specification is not an attribution
to a learner. \\

FR4 & CQ5--CQ7 &
Retrieve atomic/composite structure, direct composition and specialization,
and recorded integration or scope descriptions. Check cardinalities,
cycles, description adequacy, and the distinction from contextual
association separately. \\

FR5 & CQ9--CQ17, CQ33, CQ36 &
Retrieve qualified instructional associations and planned conditions.
Check the relevance of qualifications and preserve the distinction
between specification-level tasks and instructional units. \\

FR6 & CQ17 &
Retrieve instructional statements even when no activation is asserted.
Compare their content with the source task to assess whether the
recorded actions and deliverables are complete. \\

FR7 & CQ35, CQ36 &
Retrieve reuse of the same competence and variation in its activation
qualifications. Verify that the applicable specification remains
identifiable and is not silently redefined across contexts. \\

FR8 & CQ34 &
Retrieve alignment endpoints, external identifiers, correspondence type,
and available justification. Examine source provenance, partial scope,
and explanatory adequacy; alignment does not establish equivalence or
compliance by itself. \\

FR9 & No dedicated primary CQ &
Use complementary documentary and representational checks for expected
actions, products, observations, assessment criteria, and component
integration. Relevant instructional or integration descriptions provide
supporting material, not sufficient verification of this requirement. \\

FR10 & CQ18--CQ21, CQ24, CQ41 &
Retrieve evidence support, assessment provenance, and participation.
Check the distinction between specification and occurrence, and between
shared material and each learner's evidential basis. \\

FR11 & CQ8, CQ18, CQ21--CQ25, CQ28, CQ38, CQ41 &
Retrieve individual judgments, targets, evidence, participation, and
recorded dates. Verify applicable criteria and agreement among learner,
profile, constituent, and assessment-occasion paths separately. \\

FR12 & CQ26--CQ31, CQ37, CQ41 &
Retrieve proficiency attributions, schemes, and numeric interpretation
structures at the appropriate judgment granularity. Verify interval
rules or qualitative descriptors separately. CQ32a duplicates CQ26;
it does not provide aggregation. \\

FR13 & CQ25, CQ37, CQ38, CQ41 &
Retrieve learner-related records while preserving their assessment
occasions. Chronological ordering does not establish progression,
causation, or an overall learner classification. \\

FR14 & All operational families within their declared scope &
Record the executed query version, parameters, dataset, inference
configuration, expected answer, and observed result. Equivalent views
are not independent coverage; CQ39b is an optional update and CQ40d
is disabled by default. \\

\end{longtable}
\endgroup

The table uses family-level identifiers to remain consistent with the
chapter. For example, CQ20 includes evidence-origin and occurrence-context
views, while CQ24 includes participation and derived-record views. The
associated files must be identified when reporting a particular execution.

\section{Supplementary Operations and Equivalent Views}

CQ32, CQ39, CQ40, and CQ42 retain supplementary scope relative to the
mandatory functional baseline. Their interpretation follows the
implemented operations rather than the broader ambitions of the initial
question formulations:

\begin{itemize}
    \item \textbf{CQ32: existing overall attribution.}
    CQ32a retrieves a level already assigned to a whole-competence
    judgment and is equivalent to CQ26. The originally envisaged
    synthesis procedure is not implemented. Neither query aggregates
    constituent judgments.

    \item \textbf{CQ39: comparison candidates and recorded needs.}
    CQ39a identifies eligible judgments below a selected curricular
    minimum; CQ39b optionally materializes candidates; CQ39c retrieves
    explicitly recorded learning needs. Retrieving a candidate,
    recording it, and confirming an educational diagnosis remain
    distinct operations.

    \item \textbf{CQ40: missing records and below-minimum judgments.}
    CQ40a concerns mandatory competences without a judgment, CQ40b
    repeats the comparison in CQ39a, and CQ40c concerns missing
    constituent judgments. CQ40d remains disabled because applying a
    whole-competence minimum to every constituent requires an additional
    justified policy.

    \item \textbf{CQ42: documented bases and curricular discrepancies.}
    CQ42a retrieves competence-targeted needs and available comparison
    references. CQ42b repeats CQ39a, whereas CQ42c retrieves curricular
    inclusions without a judgment, preserving their mandatory or
    non-mandatory status.
\end{itemize}

Under identical parameters, dataset, and reference time, CQ39a, CQ40b,
and CQ42b implement the same comparison and should return equivalent
answers. This equivalence supports a regression check; it does not
constitute three independent demonstrations of functional coverage.
Likewise, CQ32a does not add synthesis evidence to the retrieval already
provided by CQ26.

The supplementary comparisons use an explicitly identified common
scheme and the documented temporal and ordinal eligibility conditions.
A curricular coverage minimum is not automatically a learner passing
threshold. Missing judgments, recorded proficiency below a minimum,
and stored learning-need interpretations therefore remain distinct
outcomes.

\section{Using the Matrix in Evaluation}

The mapping supports the planning and interpretation of evaluation,
while requirement satisfaction is reported through the criteria defined
in Section~\ref{subsec:ontoksd-requirement-verification} and the checks
in Table~\ref{tab:ontoksd-requirement-checks}. Conclusions must distinguish
support, partial support, lack of support, and absence of evaluation.
An unexecuted check establishes absence of evaluation, not failure of
the corresponding capability.

An empty result must be assessed against the expected behavior of the
selected case. For example, the absence of prerequisite activations in
the CSP data does not refute CQ11, and an existing judgment excludes a
competence from a query asking for missing judgments even when its level
is below a minimum. Synthetic positive and control cases can test these
behaviors, but must remain identifiable and cannot be reported as
observed CSP outcomes.

FR6 now has a direct retrieval path through CQ17, including instructional
units without competence activation. Its remaining completeness check
concerns correspondence with the original task description. FR9 still
requires complementary examination of expectations and criteria; the
criteria-related part of FR11 also extends beyond provenance retrieval.
These limitations must not be replaced by a blanket claim of complete
coverage.

The same distinction applies to proficiency: numeric structure retrieval
does not verify interval consistency or classify a score, and retrieval
of a qualitative attribution does not validate its rubric. Similarly,
retrieving multiple provenance paths does not enforce agreement among
them. Such checks complement the mapped queries rather than follow
from a non-empty result.

The CQ-to-FR and FR-to-CQ views must be maintained together. A change
in a query's scope, an added operational view, or a revised requirement
mapping must be reflected in both tables and in the repository catalog.
This appendix describes the evaluated mapping; it is not a reconstruction
of the historical order in which requirements and queries were developed.
